%---------------------------------------------------------------------------
\chapter{むだ時間系の制御}
%---------------------------------------------------------------------------
本章では，むだ時間系を制御するための簡単で実用的な方法を述べるにとど
め，詳しい解析などは行わない．また，本章で扱うシステムは原則としてス
カラ系である．制御系を設計する時の基本であるが，性能の良い制御を行なう
ためには，対象のプラントの構造及びパラメータをできるだけ正確に把握する
ことが好ましい．特にむだ時間系を考える場合，むだ時間の長さを正確に知る
必要がある．ここでは，モデルが比較的正確に分かっている場合の設計につい
て述べる．
%---------------------------------------------------------------------------
\section{Smithの補償器とIMC構造}
%.............................................................................
いま，プラントの伝達関数がむだ時間要素$e^{-sL}$とむだ時間のない伝達関数
$P^*(s)$の積として，
\begin{equation}
 P(s) = P^*(s)e^{-sL}
\end{equation}
で表されるとする．ただし，$L$はむだ時間であり，$P^*(s)$は安定であると
する．このとき制御系の構成をFig.\ref{fig_smith}のように構成すること
を考える．
\begin{figure}[htbp]
 \begin{center}
  \input{figure/smith.pstex_t} % scale = 100%
  \caption{Smithの補償器}
  \label{fig_smith}
 \end{center}
\end{figure}

この図で，2重のボックスで示された部分はSmithの予測器と呼ばれる．これは，
\begin{equation}
 z = Pu - (P - P^*)u = P^* u = y e^{sL}
\end{equation}
となり，$z$が出力$y$の$L$秒先の値を予測しているからである．この時，$r$
から$y$までの伝達関数は，
\begin{equation}
 \frac{Y}{R} = \frac{CP}{1 + CP^*} = \frac{C P^*}{1 + CP^*}e^{-sL}
\end{equation}
となり，ちょうどむだ時間がない場合の閉ループ伝達関数にむだ時間要素の積
という形になる．これより，コントローラ$C$の設計は，むだ時間とは独立に
従来の方法を用いることができることが分かる．ただし，この構造で注意しな
ければならないことは，$P$が安定でなければならないことである．実際，
Fig.\ref{fig_smith}を等価変換するとFig.\ref{fig_smith_eqv}を得る．
\begin{figure}[htbp]
 \begin{center}
  \input{figure/smithequiv.pstex_t} % scale = 100 %
  \caption{Smithの補償器と等価なシステム}
  \label{fig_smith_eqv}
 \end{center}
\end{figure}

もし，システム$P$が不安定である場合，コントローラ側の$P$の出力と実際の
出力がキャンセルされてしまい$e$には反映されず，全体を安定化するこ
とはできない．Fig.\ref{fig_smith_eqv}において，$C$と$P^*$の部分を一つの
コントローラ$K$として
考えた時，このような制御系をIMC(Internal Model Control)構造という．
(Fig.\ref{fig_imc})
\begin{figure}[htbp]
 \begin{center}
  \input{figure/imc.pstex_t} % scale = 100%
  \caption{IMC構造}
  \label{fig_imc}
 \end{center}
\end{figure}

この構成を用いる場合，ロバスト性をあまり考慮しなくて良い場合には，$K$
として$P^{-1}$にできるだけ近いものを選んだ方が追従特性は良くなることは明
らかであろう．さらにIMCの特徴として，コントローラ内部のプラントのモデ
ルを$\bar{P}$としてFig.\ref{fig_imc_cha}のような構成にしたとき，
閉ループ系が安定で，
\begin{equation}
 K(0)\bar{P}(0) = 1
\end{equation}
の条件を満たせばシステムは1型のサーボ系となっている．実際，
\begin{equation}
 y = \frac{KP}{1 + K(P - \bar{P})} r
\end{equation}
であるから，$r = c/s$とすると最終値定理より，
\begin{eqnarray}
 y(\infty) &=& \lim_{s \rightarrow 0}s 
  \frac{KP}{1 + K(P - \bar{P})}\frac{c}{s} \nonumber \\
 &=& \frac{K(0)P(0)}{1 - K(0)\bar{P}(0) + K(0)P(0)}c = c
\end{eqnarray}
となり出力は目標値に一致する．また，システムの入力部に飽和などの非線形
要素がある場合でも，$\bar{P}$の前に同じ非線形要素を挿入してコントロー
ラを設計することでシステムをロバスト化することができる．
\begin{figure}[htbp]
 \begin{center}
  \input{figure/imc_cha.pstex_t} % scale = 100%
  \caption{IMCの特徴}
  \label{fig_imc_cha}
 \end{center}
\end{figure}
%---------------------------------------------------------------------------
\section{むだ時間を用いてシステムを離散化する方法}
%---------------------------------------------------------------------------
前節と同じ伝達関数で与えられるプラントを考え，このシステムにゼロ次ホー
ルダを付けて制御することを考える．ただし，プラントは安定でなくとも良く，
$T$をサンプリングインターバルとして$L = nT$が成り立つように$T$を選ぶもの
とする．この時，ゼロ次ホールダを付けたシステムの伝達関数は
Fig.\ref{fig_delay1}のように考えることができる．

\begin{figure}[htbp]
 \begin{center}
  \input{figure/delay1.pstex_t} % scale = 100%
  \caption{むだ時間要素の分解1}
  \label{fig_delay1}
 \end{center}
\end{figure}

ここで入力を$u$，各むだ時間要素の出力を$u_i$とし，$P^*(z)$の差分状態空
間表現を
\begin{eqnarray*}
 x((k+1)T) &=& \Phi x(kT) + \Gamma u(kT) \\
 y(kT)     &=& C x(kT) + D u(kT)
\end{eqnarray*}
とすれば，$u_i$を状態と考えてつぎの拡大系を得る．
\begin{eqnarray}
 \left(\begin{array}{c}
        x \\
	u_n \\
	\vdots \\
	u_1
       \end{array}\right)
 ((k+1)T) &=&
 \left(\begin{array}{ccccc}
        \Phi & \Gamma &   &        & \\
	0    & 0      & 1 &        & \\
	     &        &   & \ddots & \\
	     &        &   &        & 1 \\
	     &        &   &        & 0
	\end{array}\right)
 \left(\begin{array}{c}
        x \\
	u_n \\
	\vdots \\
	u_1
	\end{array}\right)
 u(kT) +
 \left(\begin{array}{c}
        0 \\
	0 \\
	\vdots \\
	1
	\end{array}\right)
 u(kT) \\
 \left(\begin{array}{c}
        y \\
	u_n \\
	\vdots \\
	u_1
	\end{array}\right)
 (kT) &=&
 \left(\begin{array}{ccccc}
        C & D & 0 & \cdots & 0 \\
	0 &   &   &   I    & 
	\end{array}\right)
 \left(\begin{array}{c}
        x \\
	u_n \\
	\vdots \\
	u_1
       \end{array}\right)
 (kT)
\end{eqnarray}

また，出力側にむだ時間が存在する場合にはFig.\ref{fig_delay2}のように分解して
次の拡大系を得る．
\begin{figure}[htbp]
 \begin{center}
  \input{figure/delay2.pstex_t} % scale = 100%
  \caption{むだ時間要素の分解2}
  \label{fig_delay2}
 \end{center}
\end{figure}
\begin{eqnarray}
 \left(\begin{array}{c}
        y_n \\
	\vdots \\
	y_1 \\
	x
       \end{array}\right)
 ((k+1)T) &=&
 \left(\begin{array}{ccccc}
        0 & 1 &        &   & \\
	  & 0 & \ddots &   & \\
	  &   &        & 1 & \\
	  &   &        & 0 & C \\
	  &   & 0      &   & \Phi
	\end{array}\right)
 \left(\begin{array}{c}
        y_n \\
	\vdots \\
	y_1 \\
	x
       \end{array}\right)
 (kT) +
 \left(\begin{array}{c}
        0 \\
	\vdots \\
	0 \\
	D \\
	\Gamma
	\end{array}\right) u(kT) \\
 y_n(kT) &=&
  \left(\begin{array}{ccc}
         1 & 0 & \cdots
	\end{array}\right)
  \left(\begin{array}{c}
         y_n \\
	 \vdots \\
	 y_1 \\
	 x
	\end{array}\right)(kT)
\end{eqnarray}
従って，制御系はこの状態空間表現に基づいて設計すれば良い．
